\section{Introduction}
As Cyber-Physical Systems (CPS) become increasingly embedded within critical infrastructure, the tight coupling between their computational and physical layers poses significant challenges for formal verification and safety assurance. CPS have gained popularity both in industry and the research community and are represented by many varied mission critical applications~\cite{rajkumar2010cyber}. Discrete-time dynamical systems are important mathematical models used to describe the characteristics of state changes in control systems. While these models can effectively capture system behavior, their infinite state space presents challenges for verification. Current methods for verifying temporal properties of these systems include the state-triplet method~\cite{wongpiromsarn2015automata} and closure certificates~\cite{murali2024closure}. However, these methods introduce new concerns: (i) The state-triplet method suffers from conservatism; (ii) The closure certificates have low synthesis efficiency. Developing a more efficient approach that reduces conservatism for verifying temporal properties in these systems remains a significant challenge.

Prior research has focused on formal verification of temporal properties---including safety, reachability, and reach-avoid specifications---within dynamical systems \cite{xue2021reach, vzikelic2023learning, ansaripour2023learning, vzikelic2023compositional, neustroev2025neural}. However, specifications for complex systems often require LTL~\cite{pnueli1977temporal} to capture temporal behaviors. Verifying LTL properties in control systems introduces several challenges. First, these properties describe behaviors across infinite execution trajectories~\cite{baier2008principles}, typically modeled via $\omega$-automata. When employing proof certificates for verification, the state-triplet method requires unrolling all simple paths and blocking each individually. It is non-trivial to determine the number of certificates needed for synthesis and this method is conservative. Alternatively, closure certificates face the issue of enormous search spaces, as certificates are defined on $\mathcal{X} \times Q \times \mathcal{X} \times Q$. Finally, the infinite state space and nonlinear dynamics of control systems pose major obstacles for traditional verification methods.

In this work, we consider the formal verification of Linear Temporal Logic (LTL) properties for discrete-time dynamical systems using proof certificates. As abstraction-free methods, proof certificates have gained popularity in verification and static analysis domains~\cite{chakarov2013probabilistic,chatterjee2016algorithmic,chatterjee2024sound,chatterjee2024equivalence}. Murali et al.~\cite{murali2024closure} proposed closure certificates, proof certificates defined as real-valued functions over state pairs of dynamical systems, for persistence verification, which can be used to validate systems against LTL properties. Wongpiromsarn et al.~\cite{wongpiromsarn2015automata} introduced the state-triplet method, which demonstrates that a system satisfies LTL properties by blocking all reachable paths. Inspired by these two methods, we propose transition safety certificates and transition persistence certificates. Their combined use enables verification of LTL specifications expressed via transition-based Nondeterministic Büchi Automata (NBA)~\cite{buchi1966symposium}. Based on these certificates, we present an automated synthesis method using Satisfiability Modulo Theories (SMT) solvers. We show that combining these certificates reduces the conservatism of state-triplet methods while requiring smaller search spaces than closure certificates. We validate the effectiveness of our synthesis approach through two case studies.

\textbf{Contributions.} The main contributions of this paper are:
\begin{itemize}
  %\item We propose \emph{transition certificates} for verifying LTL properties of discrete-time systems with continuous state spaces. Transition certificates reduce the search space compared to closure certificates, and they also reduce conservatism compared to the state-triplet method.
  \item We propose transition certificates, a novel notion for verifying LTL properties of discrete-time dynamical systems with continuous state spaces. This method effectively addresses the limitations of existing approaches: closure certificates, which suffer from enormous search spaces (defined on \mathcal{X} \times Q \times \mathcal{X} \times Q) leading to high computational costs, and the state-triplet method, which introduces conservatism. Our approach reduces the search space to at most $\mathcal{X} \times Q \times Q$ while mitigating conservatism.
  
  % \item We present automated synthesis methods for transition certificates based on both polynomial and neural network templates, which offer​ a balance between expressiveness and computational efficiency for different system classes.
  \item We introduce automated synthesis methods​ for transition certificates based on polynomial and neural network templates. These templates provide a balance between expressiveness and computational efficiency, enabling efficient certificate generation for diverse system classes, from linear to nonlinear dynamics.
  
  % \item Through case studies, we demonstrate the verification of dynamical systems against LTL specifications using transition certificates. Comparisons with closure certificates and neural closure certificates show that our approach achieves a significant speedup due to​ its reduced search space.
  \item We demonstrate the effectiveness​ of our certificates through case studies, including the Kuramoto oscillator and a two-room temperature model. Experimental results show that transition certificates achieve significant speedups, highlighting their practicality for LTL specifications.
\end{itemize}

The organization of this paper is as follows: Section 2 covers preliminary knowledge. Section 3 introduces our proof certificates for LTL verification and compares with other methods. Section 4 presents synthesis methods. Section 5 demonstrates the effectiveness through case studies. Section 6 concludes the paper. The proofs of the relevant lemmas are provided in the Appendix~\ref{appendix-proof}.


\textbf{Related work: } 
Verification of temporal logic properties via proof certificates has seen significant advancements. To validate discrete-time systems against LTL specifications, Wongpiromsarn et al.~\cite{wongpiromsarn2015automata} developed a state-triplet method that uses barrier certificates to block all simple paths from initial states to accepting states in the automaton, requiring enumeration of state triplets $(q_m, q_m', q_m'')$ along these paths. Murali et al.~\cite{murali2024closure} introduced closure certificates for persistence verification, which establish disjunctively well-founded transition invariants \cite{podelski2004transition} for dynamical systems. Additionally, substantial research has addressed the verification of specific temporal behaviors such as reachability, safety, and avoidance~\cite{chakarov2013probabilistic,chatterjee2016algorithmic,chatterjee2024sound,chatterjee2024equivalence}. Some adopt neural networks as templates~\cite{wang2025verifying, nadali2024neural, zhao2020synthesizing} to synthesize certificates for verifying temporal properties.


The abstraction-based approaches~\cite{vasile2016control, leahy2019control, kress2009temporal, coogan2015traffic, fainekos2009temporal, ding2014optimal, wongpiromsarn2009receding} lift the CPS model from continuous domain to discrete domain by computing a finite abstraction, e.g., finite transition systems and Markov decision processes (MDPs). We note that abstraction-based techniques have been used in the verification and synthesis of continuous-space dynamical systems against LTL properties such as the results in~\cite{henzinger1997hytech, khaled2021omegathreads, lahijanian2011temporal, rungger2016scots}. These results rely on building a finite state abstraction and then making use of model checking and synthesis techniques~\cite{baier2008principles, tabuada2009verification}. In general, the abstraction-based approaches are computationally demanding and suffer from the curse of dimensionality. Abstraction-free methods, such as proof certificates, can alleviate the computational burden to some extent.


In the context of our discussion, we assume that the NBA for an LTL specification $\phi$ has already been constructed. The translation from an LTL formula to an NBA involves an exponential complexity with respect to the size of the formula~\cite{vardi2005automata}. The complexity of constructing the complement of an NBA is known to be EXPTIME~\cite{complementnba}. Converting LTL specifications into Transition-based Generalized Büchi Automata (TGBA) typically yields smaller automata compared to Labeled Generalized Büchi Automata (LGBA)~\cite{couvreur1999fly}. The degeneralization process~\cite{giannakopoulou2002states} can convert a TGBA with multiple acceptance conditions into a Transition-based Büchi Automaton (TBA) featuring a single acceptance condition. Various tools exist for generating different automaton types from LTL formulas, such as Owl~\cite{owl} and SPOT~\cite{spot}.
